\section{The decreasing-prime tree}\label{sec:tree}

Fix the residue classes in Theorem~\ref{thm:survivors}, and use
\eqref{eq:parameters}. Set
\[
 a_\star=\frac1{100},\qquad x(p)=\frac{\log p}{\log w}.
\]
Thus $x(p)$ measures a prime's exponent in base $w$, and
$1<x(p)\le B$ for $p\in(w,z]$.
A node is a finite tuple of strictly decreasing primes in $(w,z]$.
Its product is denoted by $d$, with $d=1$ for the empty tuple. Put
\begin{equation}\label{eq:tree-state}
 J_d=\frac Yd,\qquad
 r_d=\log_wJ_d-a_\star+2,\qquad \mu_d=J_dV_0.
\end{equation}
At the root, $r_0=2B-a_\star+o(1)$ and the cutoff is $b=B$.
At a nonempty node, the cutoff is $b=x(p_{\mathrm{last}})$, and the
available primes satisfy $w<p<p_{\mathrm{last}}$. The root cutoff
includes $p=z$ when $z$ is prime. All subsequent cutoffs are strict.

Let $N_d$ count integers $1\le n\le Y$ that hit every prescribed
class at the prime factors of $d$ and avoid every prescribed class at
the primes at most $w$. Let $S_d(b)$ impose, in addition, avoidance at
all available primes above $w$ up to the cutoff. In particular, the
survivor count in Theorem~\ref{thm:survivors} is $S_1(B)$.
The simultaneous hits give one class modulo $d$, so
Lemma~\ref{lem:small-sieve} applies in its progression coordinates.

Partition the integers counted by $N_d$ but not by $S_d(b)$ according
to their smallest bad available prime. This gives the exact identity
\begin{equation}\label{eq:tree-identity}
 S_d(b)=N_d-\sum_{p\text{ available}}S_{dp}(x(p)).
\end{equation}
Although the partition uses the smallest bad prime, iteration produces
decreasing prime tuples because the child count avoids all smaller
available primes.

Nodes of even length are called lower nodes, and nodes of odd length
are called upper nodes. From a lower node include every available
child. From an upper node include a child with exponent $x=x(p)$
exactly when
\begin{equation}\label{eq:admission}
 r-x\ge \mathcal H(x),\qquad \mathcal H(x)=\max(2x,x+2),
\end{equation}
where $r$ is the parent's gap. Other children are omitted.

\begin{lemma}\label{lem:tree-lower-bound}
The alternating sum of the $N_d$ over all included nodes is a lower
bound for $S_1(B)$. More generally, one may stop at any prefix-free
collection of included lower nodes, use their exact counts $S_d(b)$,
and expand only the other nodes. The resulting expression is still
a lower bound for $S_1(B)$.
Every included node satisfies $J_d\ge w^{a_\star}$.
\end{lemma}

\begin{proof}
The tree is finite because all factors are distinct primes at most
$z$. Work upward from its leaves. At a lower node, replacing each
child by an upper bound in \eqref{eq:tree-identity} gives a lower
bound. At an upper node, omitting nonnegative child counts and
replacing the retained children by lower bounds gives an upper bound.
This proves the sign assertion by induction. An exact count at a
lower node is an admissible lower bound there; prefix-freeness ensures
that its descendants are not expanded as well.

At a nonroot lower node, admission gives
$r\ge2b$ and $r\ge b+2$. Every child exponent satisfies $1<x\le b$,
so its gap is at least 2. At the root,
$r_0-B=B-a_\star+o(1)>2$ for large $z$, giving the same conclusion
for every first child. Included lower nodes also have gap exceeding
2 by their admission rule. Hence all included nodes have
$\log_wJ_d=r_d-2+a_\star\ge a_\star$.
\end{proof}

Figure~\ref{fig:admission-and-stopping} separates the admission decision
at an upper node from the later choice of lower stops.
\begin{figure}[tbp]
\centering
\begingroup
\setlength{\unitlength}{1mm}
\setlength{\fboxsep}{2mm}
\small
\begin{picture}(154,98)
  \put(77,89){\makebox(0,0){\fbox{\parbox{58mm}{\centering
    Expanded lower node $d$\\[-1pt]
    every available child is included}}}}
  \put(77,79){\vector(0,-1){15}}
  \put(81,71){\makebox(0,0)[l]{$p$}}
  \put(77,54){\makebox(0,0){\fbox{\parbox{58mm}{\centering
    Included upper node $dp$\\[-1pt]
    test each prospective lower child}}}}
  \put(69,45){\vector(-4,-1){43}}
  \put(77,44){\vector(0,-1){12}}
  \put(85,45){\vector(4,-1){43}}
  \put(46,37){\makebox(0,0){$q$}}
  \put(81,38){\makebox(0,0)[l]{$q'$}}
  \put(110,37){\makebox(0,0){$u$}}
  \put(26,19){\makebox(0,0){\fbox{\parbox{42mm}{\centering
    Included lower node $dpq$\\[-1pt]
    chosen stop: retain\\[-1pt]
    $S_{dpq}(x(q))$ exactly}}}}
  \put(77,19){\makebox(0,0){\fbox{\parbox{40mm}{\centering
    Included lower node $dpq'$\\[-1pt]
    continue its\\[-1pt]
    lower-bound expansion}}}}
  \put(105,9.5){\dashbox{1}(46,19){\parbox{40mm}{\centering
    Prospective child $dpu$\\[-1pt]
    omitted by the\\[-1pt]
    admission rule}}}
  \put(51,3){\makebox(0,0){$q,q'$ satisfy \eqref{eq:admission}}}
  \put(128,3){\makebox(0,0){$u$ fails \eqref{eq:admission}}}
\end{picture}
\endgroup
\caption{Admission and stopping have different roles. Only selected branches
are drawn, with $q,q',u<p$. At the upper node, admission is the test
$r_{dp}-x(v)\ge\mathcal H(x(v))$ for the prospective prime $v$.
The omitted child $dpu$ does not belong to the included tree.
The admitted child $dpq$ is shown as a member of a separately chosen,
prefix-free set of lower stops; its descendants are replaced by its
exact survivor count. Other included lower nodes continue to expand.
The specific stop rule is given in Section~\ref{sec:stops}.}
\label{fig:admission-and-stopping}
\end{figure}


In particular, uniformly at included nodes,
\begin{equation}\label{eq:tree-small-counts}
 \frac{N_d}{\mu_d}\ll_{a_\star}1,
 \qquad
 \left|\frac{N_d}{\mu_d}-1\right|
 \ll e^{-c_1\log_wJ_d}
 \quad\text{when }\log_wJ_d\ge v_0.
\end{equation}
The absolute upper bound remains valid on the slightly enlarged
progressions with $J_d\ge w^{a_\star/2}$ that will arise in boxes.

\subsection{The reference tree}

Replace $N_d$ by $\mu_d$ in the unstopped tree. Divide a subtree value
by its own $\mu_d$, and denote the result by $P_\ee(r,b)$ or
$P_\oo(r,b)$ according to parity. Dependence on $w$ and on the
endpoint convention is left implicit. Since $\mu_{dp}=\mu_d/p$,
the recursions are
\begin{equation}\label{eq:reference-tree}
\begin{split}
 P_\ee(r,b)&=1-\sum_{p\text{ available}}\frac1p
                         P_\oo(r-x(p),x(p)),\\
 P_\oo(r,b)&=1-
 \sum_{\substack{p\text{ available}\\r-x(p)\ge \mathcal H(x(p))}}
 \frac1p P_\ee(r-x(p),x(p)).
\end{split}
\end{equation}
The reference value at the root is $YV_0P_\ee(r_0,B)$.
Auxiliary starting states used below follow the same recursions on
their declared domains; no actual forbidden classes are needed to
define them.

A path's \emph{harmonic weight} is the reciprocal product of its
chosen primes. Dropping all admission restrictions shows that the
sum of the weights of all decreasing-prime prefixes is
\begin{equation}\label{eq:total-prefix-mass}
 \prod_{w<p\le z}\left(1+\frac1p\right)\ll B.
\end{equation}
Indeed $\log(1+1/p)=1/p+O(p^{-2})$, and the reciprocal-prime sum is
$\log B+O(1)$.

We also use \emph{standard} paths: these keep only the admission
condition $r-x\ge2x$ at an upper node. They contain all paths of
the original tree and agree with it whenever the next exponent is
at least 2. The continuous harmonic model replaces reciprocal-prime
sums by $dx/x$, with decreasing real exponents. It will be studied
first, then compared with the actual prime paths.

\subsection{Comparison after stopping}

For later use, the comparison identity can be recorded explicitly.
Let $\mathcal T$ be the retained expanded nodes of a stopped tree,
excluding the stop nodes, and let $\mathcal S$ be its prefix-free
set of lower stops. If $L_{\mathrm{stop}}$ is its lower-bound value,
then
\begin{equation}\label{eq:stopped-comparison-identity}
\begin{split}
 L_{\mathrm{stop}}-YV_0P_\ee(r_0,B)
 ={}&\sum_{d\in\mathcal T}(-1)^{\omega(d)}(N_d-\mu_d)\\
 &+\sum_{d\in\mathcal S}
 \bigl(S_d(b_d)-\mu_dP_\ee(r_d,b_d)\bigr).
\end{split}
\end{equation}
Here a node is identified by its decreasing tuple, equivalently by
its squarefree product. To verify the identity, expand the full
reference tree down to the same stops; at each stop its remaining
subtree is exactly $\mu_dP_\ee(r_d,b_d)$. Thus a positive reference
margin, a small sum of retained absolute errors, and an aggregate
nonnegative stopped difference will prove
Theorem~\ref{thm:survivors}.
